% ============================================================
%                 CHAPTER 4: METHODOLOGY
% ============================================================
\chapter{Methodology}

\section{Core Integration Framework Strategy}

The methodology for bridging SUMO and Unity revolves primarily around the concept of "Inversion of Traffic Control." In a standard co-simulation, SUMO operates as the absolute dictator of all entities. However, for a driving simulator scenario, the control authority over the Ego Vehicle must be explicitly revoked from SUMO and handed over to Unity’s Rigidbody physics engine. Meanwhile, SUMO must retain control over the remaining background flora consisting of hundreds of autonomous agents. 

This creates a split-authority architecture that necessitates rigorous, bidirectional polling at incredibly high frequencies. If SUMO assumes the ego vehicle is executing a Krauss braking maneuver, but the human user actually accelerates inside the simulator, the two engines will become violently desynchronized. Therefore, the ego vehicle's $X, Y, Z$ and $\theta$ parameters must be forcefully overwritten in SUMO’s memory cache at the exact end of every Unity physics step loop.

\begin{figure}[htbp]
\centering
\begin{tikzpicture}[
    outerbox/.style={
        rectangle, rounded corners=6pt, draw=PrimaryNavy, line width=1.5pt,
        fill=white, minimum width=6.5cm, minimum height=4.5cm
    },
    innerbox/.style={
        rectangle, rounded corners=4pt, draw=BorderColor, line width=1pt,
        fill=SoftBg, minimum width=5.5cm, minimum height=0.7cm, align=center,
        font=\small\sffamily\color{DarkText}
    },
    commbus/.style={
        rectangle, rounded corners=6pt, draw=AccentTeal, line width=1.5pt,
        fill=AccentTeal!10, minimum width=14.5cm, minimum height=1.5cm, align=center,
        font=\bfseries\sffamily\color{PrimaryNavy}
    },
    arr/.style={
        <{Stealth[length=3mm, width=2.5mm]}-{Stealth[length=3mm, width=2.5mm]}, line width=1.5pt, color=PrimaryNavy
    }
]

    % SUMO BOX
    \node[outerbox] (sumobox) {};
    \node[anchor=north, font=\Large\bfseries\sffamily\color{PrimaryNavy}, yshift=-10pt] (sumotitle) at (sumobox.north) {SUMO Core};
    \draw[AccentTeal, line width=2pt] ([yshift=-35pt, xshift=-2.6cm]sumobox.north) -- ([yshift=-35pt, xshift=2.6cm]sumobox.north);
    
    \node[innerbox, below=0.4cm of sumotitle] (s1) {Vehicle Kinematics (Krauss/IDM)};
    \node[innerbox, below=0.15cm of s1] (s2) {Pedestrian Models (SFM)};
    \node[innerbox, below=0.15cm of s2] (s3) {TraCI TCP Server};

    % UNITY BOX
    \node[outerbox, right=1.5cm of sumobox] (unitybox) {};
    \node[anchor=north, font=\Large\bfseries\sffamily\color{PrimaryNavy}, yshift=-10pt] (unitytitle) at (unitybox.north) {Unity3D Engine};
    \draw[AccentTeal, line width=2pt] ([yshift=-35pt, xshift=-2.6cm]unitybox.north) -- ([yshift=-35pt, xshift=2.6cm]unitybox.north);
    
    \node[innerbox, below=0.4cm of unitytitle] (u1) {Universal Render Pipeline (URP)};
    \node[innerbox, below=0.15cm of u1] (u2) {NVIDIA PhysX Rigidbodies};
    \node[innerbox, below=0.15cm of u2] (u3) {Input Configuration / VRTK};

    % Communication BUS
    \node[commbus, below=1.5cm of $(sumobox.south)!0.5!(unitybox.south)$] (comm) {
        ZERO-MQ NETWORKING PROTOCOL \\[4pt]
        \normalfont\small\sffamily NetMQ C\# Wrapper \quad \textbf{|} \quad REQ-REP Handshake \quad \textbf{|} \quad JSON Serialization \quad \textbf{|} \quad $< 5$ms Latency
    };
    
    % Arrows
    \draw[arr] (sumobox.south) -- (sumobox.south |- comm.north);
    \draw[arr] (unitybox.south) -- (unitybox.south |- comm.north);
    \draw[arr, AccentTeal] (sumobox.east) -- node[above, font=\sffamily\small\color{PrimaryNavy}\bfseries] {High-Frequency Sync} (unitybox.west);

\end{tikzpicture}
\caption{System architecture detailing the separation of concerns and the NetMQ bidirectional bridge.}
\label{fig:arch}
\end{figure}

\section{Data Exchange Using ZeroMQ (ZMQ) and NetMQ}

To achieve sub-5ms latency across the boundary of the Python daemon (SUMO TraCI) and the C\# Execution Engine (Unity), a specialized message-broker pattern is initiated using ZeroMQ (ZMQ). ZMQ functions natively outside the traditional TCP stack dependencies, drastically reducing the system interrupts traditionally generated by Windows OS socket polling.

\begin{insightbox}[The NetMQ Implementation]
NetMQ, the native C\# rewrite of the ZeroMQ library, was integrated directly into Unity's plugin architecture. This prevents the traditional garbage collection (GC) stalls associated with unmanaged memory wrappers. Operating exclusively on a BackgroundWorker thread, NetMQ prevents the network polling loop from blocking the main rendering pipeline.
\end{insightbox}

\subsection{The Request-Reply (REQ-REP) Paradigm}

The chosen communication methodology is a strict synchronous handshake via a REQ-REP pattern over \texttt{tcp://localhost:5555}. Although PUB-SUB (Publish-Subscribe) allows asynchronous firing, it natively suffers from packet dropping if the subscriber evaluates frames slower than the publisher. If even a single frame of ego-vehicle telemetry is dropped before SUMO processes it, a fatal collision could occur in the background simulation resulting in erroneous macroscopic traffic deadlocks.

\textbf{Execution Sequence:}
\begin{enumerate}
    \item \textbf{Unity Request Phase:} At the end of the \texttt{FixedUpdate()} block, Unity dumps the ego vehicle coordinates into a tiny JSON string and dispatches it via the NetMQ \texttt{RequestSocket.SendFrame()}.
    \item \textbf{SUMO Ingestion:} The \texttt{runner.py} script detects the bytecode on the ZMQ \texttt{REP} socket. It calls \texttt{traci.vehicle.moveToXY()} forcefully shifting the ego vehicle block within the network abstraction.
    \item \textbf{SUMO Stepping Phase:} TraCI evaluates \texttt{traci.simulationStep()}. All active autonomous agents calculate their distance against the newly relocated human-driven ego vehicle.
    \item \textbf{SUMO Reply Phase:} \texttt{runner.py} iterates through `traci.vehicle.getIDList()`, pulling the $X, Y, \theta$, and velocity parameters for every vehicle, packing it into a massive concatenated string, and emitting it via the ZMQ `REP` confirmation socket.
    \item \textbf{Unity Resolution Phase:} Unity’s NetMQ thread catches the reply, deserializes it, updates the target coordinate dictionary, and signals the main thread interpolator.
\end{enumerate}

\section{Serialization Overhead and JSON Formatting}

One of the methodological complexities discovered involved defining the structural payload format. Directly sending massive strings such as \texttt{"car1,10.5,22.1,90|car2,11.2,2.3,180"} via byte arrays is highly efficient but brittle (a single floating-point corruption ruins the array). Using full \texttt{JSON} structures provides safety but inflates string lengths dramatically.

The JSON methodology was strictly optimized by stripping metadata:
\begin{lstlisting}
{
  "t": 34.50,
  "v": {
    "c_01": [152.3, 44.1, 90.0, 15.5],
    "c_02": [188.1, 41.0, 90.0, 12.0]
  }
}
\end{lstlisting}

This concise structure (ID mapped to an array sequence of $[X, Y, \text{Angle}, \text{Speed}]$) limits the payload to approximately 40 bytes per vehicle. Thus, transmitting 100 background agents scales to merely 4.0 KB per frame, fitting comfortably within a single maximum transmission unit (MTU) to avoid IP fragmentation overhead natively across the loopback adapter.

\section{Scenario Design and Traffic Injection}

The final methodological step is the dynamic scenario manipulation capability. TraCI natively supports dynamic insertions. When attempting to simulate edge cases—for instance, an aggressive driver illegally merging—creating strict XML `.rou.xml` templates is tedious. 

Instead, the Python \texttt{runner.py} acts as an omnipresent hypervisor. It continuously monitors the human ego vehicle's velocity. If the ego vehicle approaches an intersection at high speed, the \texttt{runner.py} script conditionally injects a hazard vehicle algorithmically by executing \texttt{traci.vehicle.add()}, forcing a background car to run a red light right as the player approaches. 

This algorithmic generation of hazard states transforms the simulation from a passive visualizer into an aggressive, rigorous driving simulator scenario designer capable of repeatedly forcing users into identical high-stress, millimeter-perfect crash-avoidance situations without manual scripting overhead in the Unity Timeline.
